\section{Zadanie 3 (Dominik Korsa)}
Niech \(X\) będzie zmienną losową z rozkładu normalnego
\(N(0, \sigma^2)\), zaś \(Z\) zmienną przyjmującą wartości ze zbioru
\(\{-1, +1\}\) z rozkładem jednostajnym. Niech \(Y = XZ\).

\begin{enumerate}[label=(\roman*)]
  \item\label{set10.task03.1} Pokaż, że zmienna \(Y\) ma rozkład \(N(0, \sigma^2)\).
  \item\label{set10.task03.2} Wyjaśnij, dlaczego zmienne \(X, Y\) nie są niezależne.
  \item\label{set10.task03.3} Pokaż, że zmienne \(X\), \(Y\) są nieskorelowane (czyli \(Cov(X, Y) = 0\)).
\end{enumerate}

\subsection*{Rozwiązanie podpunktu \ref{set10.task03.1}}
Dystrybuanta zmiennej \(X\) to z własności rozkładu normalnego
\[
  F_X(x)
  = P(X \leq x)
  = P\left( \frac{X}{\sigma} \leq \frac{x}{\sigma} \right)
  = \Phi\left(\frac{x}{\sigma}\right).
\]
\(\frac{X}{\sigma}\) jest zmienną o standardowym rozkładzie normalnym,
a \(\Phi\) oznacza dystrybuantę standardowego rozkładu normalnego.

Pokażmy, że \(XZ\) ma taką samą dystrybuantę, a więc też taki sam rozkład, co \(X\).
\begin{align*}
  F_{XZ}(x)
  &= P(XZ \leq x)\\
  &= P\Big((Z = 1 \land X \leq x)\; \lor \; (Z = -1 \land -X \leq x) \Big)
\end{align*}
Z rozłączności zdarzeń oraz niezależności \(X\) i \(Z\):
\begin{align*}
  F_{XZ}(x)
  &= P(Z = 1)P(X \leq x) + P(Z=-1)P(-X \leq x)\\
  &= \frac{1}{2} P(X \leq x) + \frac{1}{2} P(-X \leq x)
\end{align*}

\(P(X \leq x)\) przekształciliśmy już powyżej, zróbmy to samo z \(P(-X \leq x)\).
Ostatnią równość mamy z wykładu:
\[
  P(-X \leq x)
  = P\left( \frac{X}{\sigma} \geq \frac{-x}{\sigma} \right)
  = 1-P\left( \frac{X}{\sigma} \leq \frac{-x}{\sigma} \right)
  = 1 - \Phi\left( \frac{-x}{\sigma} \right)
  = \Phi\left( \frac{x}{\sigma} \right)
\]

Zatem:
\[
  F_{XZ}(x) = \frac{1}{2} \Phi\left( \frac{x}{\sigma} \right)
  + \frac{1}{2} \Phi\left( \frac{x}{\sigma} \right)
  = \Phi\left( \frac{x}{\sigma} \right).
\]
\(X\) i \(Y = XZ\) mają taką samą dystrybuantę, więc też taki sam rozkład. 

\subsection*{Rozwiązanie podpunktu \ref{set10.task03.2}}
Załóżmy \textbf{nie wprost}, że \(X\) i \(Y\) są niezależne.
W takim razie:
\[
  P\big(X \in [-1, 1] \land Y \in [-1, 1]\big)
  = P\big(X \in [-1, 1]\big) \cdot P\big(Y \in [-1, 1]\big).
\]
Ale \(X \in [-1, 1]\) jest tym samym zdarzeniem, co \(Y \in [-1, 1]\), zatem
\[
  P\big(X \in [-1, 1] \land Y \in [-1, 1]\big)
  = P\big(X \in [-1, 1]\big)
  = P\big(Y \in [-1, 1]\big).
\]

Z poprzedniego równania otrzymujemy oczywiście fałszywy wniosek,
że \(P(X \in [-1, 1]) = 1\), zatem założenie jest sprzeczne.

\subsection*{Rozwiązanie podpunktu \ref{set10.task03.3}}
Rozpiszmy \(Cov(X, Y)\) z definicji:
\begin{align*}
  Cov(X, Y)
  &= E[X Y] - \underbrace{E[X]}_{=0} \cdot E[Y]\\
  &= E[X Y]\\
  &= P(Z = -1) E[XY \mid Z = -1] + P(Z = +1) E[XY \mid Z = +1]
\end{align*}
Rozpiszmy wiedząc, że \(Y = ZX\) i korzystając z niezależności \(X\) i \(Z\):
\begin{itemize}
  \item \(E[XY \mid Z = -1]
    = E[X (-X) \mid Z = -1]
    = E[-X^2]
    = -E[X^2]
  \)
  \item \(E[XY \mid Z = +1]
    = E[X \cdot X \mid Z = +1]
    = E[X^2]
  \)
\end{itemize}
Zatem
\[
  Cov(X, Y)
  = -\frac{1}{2} E[X^2] + \frac{1}{2} E[X^2]
  = 0.
\]
